
\section{Tối ưu hóa Nguyên-Hỗn hợp}
\label{sec:mip_background}

Tối ưu hóa nguyên-hỗn hợp là nền tảng tính toán của các kỹ thuật được giới thiệu trong đồ án này.
Nhìn chung, một \textbf{Bài toán Quy hoạch Nguyên-Hỗn hợp (Mixed-Integer Program~-- MIP)} có thể rất khó
giải và thời gian chạy của bộ giải nguyên-hỗn hợp có thể tăng rất nhanh (theo cấp số nhân) cùng kích thước
bài toán.
Tuy nhiên, đây là trường hợp xấu nhất; trên thực tế, có thể xây dựng các MIP hiệu quả cao
được giải nhanh trong hầu hết các trường hợp mà chúng ta gặp phải.
Mục tiêu của phần này là hai mặt: thứ nhất, giới thiệu các kiến thức nền tảng về tối ưu hóa
nguyên-hỗn hợp (MIP là gì, cách phân loại MIP, và các thuật toán nào có thể dùng để giải chúng);
thứ hai, giới thiệu các ký hiệu và công cụ cần thiết để phân biệt giữa các MIP hiệu quả và không hiệu quả.

\subsection{Bài toán Quy hoạch Nguyên-Hỗn hợp}

Một MIP là bài toán tối ưu hóa với các biến liên tục $x \in \mathbb{R}^n$ và các biến rời rạc
(nguyên) $y \in \mathbb{Z}^m$.
Cho trước một hàm mục tiêu $f : \mathbb{R}^{n+m} \to \mathbb{R}$ và một tập ràng buộc
$\mathcal{S} \subseteq \mathbb{R}^{n+m}$, một MIP tổng quát có thể được phát biểu như sau:
\begin{subequations}
\label{eq:mip_full}
\begin{align}
\text{minimize} \quad & f(x, y) \label{eq:mip_full_obj}\\
\text{subject to} \quad & (x, y) \in \mathcal{S}, \label{eq:mip_full_con}\\
& y \in \mathbb{Z}^m. \label{eq:mip_full_int}
\end{align}
\end{subequations}
(Ràng buộc $y \in \mathbb{Z}^m$ được phát biểu tường minh vì, nếu không có quy định khác, các biến của
bài toán tối ưu luôn được giả định là có giá trị thực.)
Chúng ta gọi tập hợp
\[
\mathcal{T} := \mathcal{S} \cap \mathbb{R}^n \times \mathbb{Z}^m
\]
là \textit{tập chấp nhận được} (feasible set) của MIP~\eqref{eq:mip_full} và các phần tử
$(x, y) \in \mathcal{T}$ là các \textit{nghiệm chấp nhận được}.
Bài toán MIP là \textit{khả thi} nếu tồn tại nghiệm chấp nhận được và \textit{bất khả thi} trong
trường hợp ngược lại.
Một nghiệm chấp nhận được $(x^{\text{opt}}, y^{\text{opt}})$ là \textit{tối ưu} nếu
$f(x^{\text{opt}}, y^{\text{opt}}) \leq f(x, y)$ với mọi $(x, y) \in \mathcal{T}$.
Nếu nghiệm tối ưu tồn tại, ta ký hiệu $f^{\text{opt}} := f(x^{\text{opt}}, y^{\text{opt}})$
là \textit{giá trị tối ưu} của MIP.

Nếu bỏ ràng buộc~\eqref{eq:mip_full_int} khỏi MIP, ta thu được bài toán tối ưu sau:
\begin{subequations}
\label{eq:mip_relax}
\begin{align}
\text{minimize} \quad & f(x, y) \\
\text{subject to} \quad & (x, y) \in \mathcal{S}.
\end{align}
\end{subequations}
Bài toán này được gọi là \textbf{bài toán nới lỏng} (relaxation) của MIP.
Ta ký hiệu nghiệm tối ưu của bài toán nới lỏng là $(x^{\text{relax}}, y^{\text{relax}})$
và $f^{\text{relax}} := f(x^{\text{relax}}, y^{\text{relax}})$ là giá trị tối ưu của nó.
Lưu ý rằng
\[
f^{\text{relax}} \leq f^{\text{opt}},
\]
do việc bỏ một ràng buộc khỏi bài toán tối ưu (tối thiểu hóa) chỉ có thể làm giảm giá trị tối ưu.
Như được thảo luận ở phần dưới, độ chặt của bất đẳng thức này đóng vai trò trung tâm trong tính hiệu
quả của một MIP.
Nói một cách ngắn gọn, một MIP có thể được giải hiệu quả nếu bài toán nới lỏng của nó có thể được giải
nhanh và khoảng cách $f^{\text{opt}} - f^{\text{relax}} \geq 0$ là nhỏ.

Các MIP được phân loại theo các tính chất của hàm $f$ và tập $\mathcal{S}$,
những yếu tố này có thể ảnh hưởng đáng kể đến khả năng giải MIP một cách hiệu quả.
Dưới đây chúng ta định nghĩa các lớp MIP quan trọng nhất trong phạm vi đồ án này.

\subsection{Bài toán Quy hoạch Nhị phân-Hỗn hợp}

Một \textbf{Bài toán Quy hoạch Nhị phân-Hỗn hợp (Mixed-Boolean Program~-- MBP)} là một MIP
trong đó các biến rời rạc chỉ có thể nhận giá trị nhị phân: $y \in \{0, 1\}^m$.
Tương đương với đó, một MBP là bài toán có dạng~\eqref{eq:mip_full} với
\[
\mathcal{S} \subset \mathbb{R}^n \times [0, 1]^m.
\]
Tất cả các MIP mà chúng ta sẽ gặp trong đồ án này thực ra đều là các MBP.
Tuy nhiên, thuật ngữ MBP không phổ biến và, mặc dù về mặt kỹ thuật các bài toán của chúng ta là MBP,
chúng ta vẫn gọi chúng là MIP cho nhất quán với quy ước của lĩnh vực.

\subsection{Bài toán Quy hoạch Lồi Nguyên-Hỗn hợp}

Nếu hàm mục tiêu $f$ và tập ràng buộc $\mathcal{S}$ đều là lồi, ta gọi bài toán~\eqref{eq:mip_full}
là \textbf{Bài toán Quy hoạch Lồi Nguyên-Hỗn hợp (Mixed-Integer Convex Program~-- MICP)}.
MICP là một lớp cơ bản của MIP vì các bài toán nới lỏng của chúng là các bài toán tối ưu hóa lồi,
vốn trong hầu hết các trường hợp có thể được giải rất nhanh.
Để nhấn mạnh giả thiết tính lồi, ta gọi bài toán nới lỏng của một MICP là \textit{nới lỏng lồi}
(convex relaxation).
Sử dụng thuật toán Nhánh-và-Cận (Branch-and-Bound~-- BB), hầu hết các MICP có thể được giải một cách
đáng tin cậy đến mức tối ưu toàn cục, mặc dù thuật toán có thể mất thời gian đáng kể.

Trái ngược với lớp MICP, chúng ta gọi \textbf{Bài toán Quy hoạch Phi-lồi Nguyên-Hỗn hợp
(Mixed-Integer Non-Convex Program~-- MINCP)} là một MIP trong đó hàm mục tiêu $f$ và/hoặc tập ràng buộc
$\mathcal{S}$ là không lồi.
Hầu hết các MINCP đều không thể giải được trong thời gian thực tế; ngoại lệ chính là các MINCP
rất nhỏ trong đó tập chấp nhận được $\mathcal{T}$ có thể được rời rạc hóa mịn và tìm kiếm vét cạn.

\subsection{Bài toán Quy hoạch Nón Nguyên-Hỗn hợp}

Nếu hàm mục tiêu $f$ là tuyến tính và tập ràng buộc $\mathcal{S}$ là một tập lồi đóng dưới dạng nón,
bài toán~\eqref{eq:mip_full} được gọi là \textbf{Bài toán Quy hoạch Nón Nguyên-Hỗn hợp
(Mixed-Integer Conic Program~-- MIKP)}.
Các lớp con của nhóm bài toán này được phân loại như sau:

\begin{itemize}
    \item \textbf{Quy hoạch Tuyến tính Nguyên-Hỗn hợp (Mixed-Integer Linear Program~-- MILP)}
    khi $\mathcal{S}$ là một đa diện (polyhedron).
    \item \textbf{Quy hoạch Nón Bậc hai Nguyên-Hỗn hợp (Mixed-Integer Second-Order Cone Program~-- MISOCP)}
    khi $\mathcal{S}$ là tập lồi bậc hai.
    \item \textbf{Quy hoạch Bán xác định Nguyên-Hỗn hợp (Mixed-Integer Semidefinite Program~-- MISDP)}
    khi $\mathcal{S}$ là một spectrahedron.
\end{itemize}

Các bài toán nới lỏng của chúng được đặt tên theo cách tương tự; ví dụ, bài toán nới lỏng của một MILP
được gọi là \textit{nới lỏng tuyến tính} (linear relaxation).

MILP đóng vai trò đặc biệt quan trọng trong tối ưu hóa nguyên-hỗn hợp, quan trọng hơn cả vai trò
của LP trong tối ưu hóa lồi, vì hai lý do.
Lý do thứ nhất mang tính hình học: các đa diện là hình dạng đơn giản nhất có thể bao quanh một tập
hữu hạn các điểm, và do đó là ứng viên tự nhiên để mô hình hóa phần rời rạc của một MIP.
Lý do thứ hai mang tính thuật toán: thuật toán đơn hình (simplex) cho tối ưu tuyến tính có thể
được tích hợp vào thủ tục BB hiệu quả hơn so với thuật toán điểm trong được dùng cho các chương trình
nón tổng quát hơn.

Ngoài ra, một \textbf{Bài toán Quy hoạch Toàn phương Nguyên-Hỗn hợp (Mixed-Integer Quadratic
Program~-- MIQP)} là một MIP có dạng~\eqref{eq:mip_full} với $f$ là hàm bậc hai và $\mathcal{S}$ là
đa diện.
Các bài toán này có thể được biểu diễn dưới dạng MISOCP, nhưng tương tự như MILP, chúng có thể được
giải hiệu quả hơn bằng các phương pháp BB chuyên biệt khai thác tính đa diện của tập ràng buộc.
Ngày nay, các MISOCP và MISDP cũng có thể được giải khá hiệu quả ngay cả với các bộ giải mã nguồn mở
(xem, ví dụ, Pajarito~\cite{lubin2018}).
